\documentclass[a4paper,11pt,twoside]{ThesisStyle}

% \usepackage[spanish]{babel} % division de silabas en spanish.
% \usepackage{ucs}
% \usepackage[utf8x]{inputenc}

\usepackage[spanish]{babel}
\usepackage[utf8]{inputenc}

% AMS Math
\usepackage{amsmath, amsthm, amssymb}
% \usepackage{amsmath,amssymb}
% \usepackage[latin1]{inputenc}

\usepackage{longtable}
\usepackage{multirow}

\usepackage{listings}
\usepackage{color}
\usepackage{textcomp}

\definecolor{listinggray}{gray}{0.9}
\definecolor{lbcolor}{rgb}{0.9,0.9,0.9}
\definecolor{gris}{rgb}{0.8,0.8,0.8}
\definecolor{azul}{rgb}{0.0,0.3,0.6}

\lstdefinelanguage{specmag}{
  keywords={compute, end, all, semantic, domain, attributes, rules, sort, op, function, infix, prefix, postfix, syn, inh, left, right, non_assoc, and, and_eq, asm, auto, bitand, bitor, break, case, catch, class, compl, const, const_cast, continue, default, delete, do, double, dynamic_cast, else, enum, explicit, export, extern, false, for, friend, goto, if, inline, long, mutable, namespace, new, not, not_eq, operator, or, or_eq, private, protected, public, register, reinterpret_cast, return, short, signed, sizeof, static, static_cast, struct, switch, template, this, throw, true, try, typedef, typeid, typename, union, unsigned, using, virtual, void, volatile, wchar_t, while, xor, xor_eq, bool, char, float, int, string, repeat, until},
  sensitive=true,
  morecomment=[s]{/*}{*/},
  morecomment=[l]{\//},
  morestring=[d]{"}
}


\lstset{
    backgroundcolor=\color{lbcolor},
    tabsize=2,
%   rulecolor=,
%   language=matlab,
%   basicstyle=\tiny,
%   basicstyle=\renewcommand{\baselinestretch}{0.97}\small\tt
    upquote=true,
%     aboveskip={1\baselineskip},
    columns=[c]fixed,
    showstringspaces=false,
    extendedchars=true,
    breaklines=true,
    prebreak = \raisebox{0ex}[0ex][0ex]{\ensuremath{\hookleftarrow}},
    frame=single,
    showtabs=false,
    showspaces=false,
    showstringspaces=false,
    identifierstyle=\ttfamily,
    keywordstyle=\color[rgb]{0.1,0.1,0.6}\bfseries,
    commentstyle=\color[rgb]{0.133,0.545,0.133},
    stringstyle=\color[rgb]{0.627,0.126,0.941},
    keepspaces=true,
    escapeinside=``,
    numbersep=5pt,
    numberstyle=\tiny
}


\usepackage[T1]{fontenc}
% \usepackage[left=3cm,right=2cm,top=4cm,bottom=3cm,includefoot,includehead,headheight=13.6pt]{geometry}
\usepackage[left=1.3in,right=1.1in,top=1.1in,bottom=1.1in,includefoot,includehead,headheight=13.6pt]{geometry}
\renewcommand{\baselinestretch}{1.05}

% Table of contents for each chapter
\usepackage[nottoc, notlof, notlot]{tocbibind}
\usepackage{minitoc}
\setcounter{minitocdepth}{2}
\mtcindent=15pt
\renewcommand\mtctitle{Contenido}
\renewcommand\mlftitle{Figuras}
\renewcommand\mlttitle{Tablas}

\usepackage{aecompl}

% Glossary / list of abbreviations

\usepackage[intoc]{nomencl}
\renewcommand{\nomname}{Lista de Abreviaturas}

\newcommand{\maggen}{\textbf{magGen}}
\newcommand{\boost}{\textit{\textbf{Boost}}}
\newcommand{\spirit}{\textit{\textbf{Spirit}}}
\newcommand{\flecha}{$\hfill \Longrightarrow \hfill$}

\makenomenclature

\usepackage{ifpdf}

\ifpdf
  \usepackage[pdftex]{graphicx}
  \DeclareGraphicsExtensions{.eps,.jpg}
  \usepackage[a4paper,pagebackref,hyperindex=true]{hyperref}
\else
  \usepackage{graphicx}
  \DeclareGraphicsExtensions{.ps,.eps}
  \usepackage[a4paper,dvipdfm,pagebackref,hyperindex=true]{hyperref}
\fi

\graphicspath{{.}{images/}}

% Links in pdf
\usepackage{color}
\definecolor{linkcol}{rgb}{0,0,0.4} 
\definecolor{citecol}{rgb}{0.5,0,0} 

% Change this to change the informations included in the pdf file

% See hyperref documentation for information on those parameters

\hypersetup
{
bookmarksopen=true,
pdftitle="\maggen",
pdfauthor="Kilmurray - Picco'", 
pdfsubject="Generador de Evaluadores estáticos para MAG", %subject of the document
%pdftoolbar=false, % toolbar hidden
pdfmenubar=true, %menubar shown
pdfhighlight=/O, %effect of clicking on a link
colorlinks=true, %couleurs sur les liens hypertextes
pdfpagemode=None, %aucun mode de page
pdfpagelayout=SinglePage, %ouverture en simple page
pdffitwindow=true, %pages ouvertes entierement dans toute la fenetre
linkcolor=linkcol, %couleur des liens hypertextes internes
citecolor=citecol, %couleur des liens pour les citations
urlcolor=linkcol %couleur des liens pour les url
}

% definitions.
% -------------------

\setcounter{secnumdepth}{3}
\setcounter{tocdepth}{2}

% Some useful commands and shortcut for maths:  partial derivative and stuff

\newcommand{\pd}[2]{\frac{\partial #1}{\partial #2}}
\def\abs{\operatorname{abs}}
\def\argmax{\operatornamewithlimits{arg\,max}}
\def\argmin{\operatornamewithlimits{arg\,min}}
\def\diag{\operatorname{Diag}}
\newcommand{\eqRef}[1]{(\ref{#1})}

\usepackage{rotating}                    % Sideways of figures & tables
%\usepackage{bibunits}
%\usepackage[sectionbib]{chapterbib}          % Cross-reference package (Natural BiB)
%\usepackage{natbib}                  % Put References at the end of each chapter
                                         % Do not put 'sectionbib' option here.
                                         % Sectionbib option in 'natbib' will do.
\usepackage{fancyhdr}                    % Fancy Header and Footer

% \usepackage{txfonts}                     % Public Times New Roman text & math font
  
%%% Fancy Header %%%%%%%%%%%%%%%%%%%%%%%%%%%%%%%%%%%%%%%%%%%%%%%%%%%%%%%%%%%%%%%%%%
% Fancy Header Style Options

\pagestyle{fancy}                       % Sets fancy header and footer
\fancyfoot{}                            % Delete current footer settings

%\renewcommand{\chaptermark}[1]{         % Lower Case Chapter marker style
%  \markboth{\chaptername\ \thechapter.\ #1}}{}} %

%\renewcommand{\sectionmark}[1]{         % Lower case Section marker style
%  \markright{\thesection.\ #1}}         %

\fancyhead[LE,RO]{\bfseries\thepage}    % Page number (boldface) in left on even
% pages and right on odd pages
\fancyhead[RE]{\bfseries\nouppercase{\leftmark}}      % Chapter in the right on even pages
\fancyhead[LO]{\bfseries\nouppercase{\rightmark}}     % Section in the left on odd pages

\let\headruleORIG\headrule
\renewcommand{\headrule}{\color{black} \headruleORIG}
\renewcommand{\headrulewidth}{1.0pt}
\usepackage{colortbl}
\arrayrulecolor{black}

\fancypagestyle{plain}{
  \fancyhead{}
  \fancyfoot{}
  \renewcommand{\headrulewidth}{0pt}
}

\usepackage[plain]{algorithm}
\usepackage[noend]{algorithmic}

%%% Clear Header %%%%%%%%%%%%%%%%%%%%%%%%%%%%%%%%%%%%%%%%%%%%%%%%%%%%%%%%%%%%%%%%%%
% Clear Header Style on the Last Empty Odd pages
\makeatletter

\def\cleardoublepage{\clearpage\if@twoside \ifodd\c@page\else%
  \hbox{}%
  \thispagestyle{empty}%              % Empty header styles
  \newpage%
  \if@twocolumn\hbox{}\newpage\fi\fi\fi}

\makeatother
 
%%%%%%%%%%%%%%%%%%%%%%%%%%%%%%%%%%%%%%%%%%%%%%%%%%%%%%%%%%%%%%%%%%%%%%%%%%%%%%% 
% Prints your review date and 'Draft Version' (From Josullvn, CS, CMU)
\newcommand{\reviewtimetoday}[2]{\special{!userdict begin
    /bop-hook{gsave 20 710 translate 45 rotate 0.8 setgray
      /Times-Roman findfont 12 scalefont setfont 0 0   moveto (#1) show
      0 -12 moveto (#2) show grestore}def end}}
% You can turn on or off this option.
% \reviewtimetoday{\today}{Draft Version}
%%%%%%%%%%%%%%%%%%%%%%%%%%%%%%%%%%%%%%%%%%%%%%%%%%%%%%%%%%%%%%%%%%%%%%%%%%%%%%% 

\newenvironment{maxime}[1]
{
\vspace*{0cm}
\hfill
\begin{minipage}{0.5\textwidth}%
%\rule[0.5ex]{\textwidth}{0.1mm}\\%
\hrulefill $\:$ {\bf #1}\\
%\vspace*{-0.25cm}
\it 
}%
{%

\hrulefill
\vspace*{0.5cm}%
\end{minipage}
}

\let\minitocORIG\minitoc
\renewcommand{\minitoc}{\minitocORIG \vspace{1.5em}}

\usepackage{slashbox}

\newenvironment{bulletList}%
{ \begin{list}%
	{$\bullet$}%
	{\setlength{\labelwidth}{25pt}%
	 \setlength{\leftmargin}{30pt}%
	 \setlength{\itemsep}{\parsep}}}%
{ \end{list} }


%\newenvironment{definition}[1][Definicion]{\begin{trivlist}
%\item[\hskip \labelsep {\bfseries #1}]}{\end{trivlist}}

\newtheorem{definition}{Definici\'on}[section]
\newtheorem{theorem}{Teorema}[section]


\newcounter{ruleAGcounter}
\newenvironment{AG}
{
 \setcounter{ruleAGcounter}{0} \begin{center} \begin{figure}[!hbtp]
 \rule{\textwidth}{.5pt} \footnotesize
}
{\normalsize \rule{\textwidth}{.5pt} \end{figure} \end{center}}

\newcommand{\newproduction}[2]{$p_\arabic{ruleAGcounter}$: 
            \emph{#1} $\rightarrow$ \emph{#2}
            \addtocounter{ruleAGcounter}{1}\\}

\newenvironment{Attribution}
{\hspace*{1.8cm}\textbf{attribution}\\}
{\hspace*{1.8cm}\textbf{end}\\}
\newcommand{\attribution}[1]{\hspace*{2cm}\emph{#1} \\}


\newcommand{\clausderiva} [1] { 
  \makebox[0pt][l]{ \raisebox{-1.2ex}{\hspace{3pt}\tiny{#1}} }
  \makebox[0pt][l]{ \raisebox{0.9ex}{\hspace{3pt}\tiny{*}} }
  \makebox[1.1\width]{$ \Longrightarrow $}
}

\renewcommand{\epsilon}{\varepsilon}

% centered page environment

\newenvironment{vcenterpage}
{\newpage\vspace*{\fill}\thispagestyle{empty}\renewcommand{\headrulewidth}{0pt}}
{\vspace*{\fill}}

\newcommand{\urllink}[1]{\htmladdnormallink{#1}{#1}}

\newenvironment{items}{
\begin{itemize}
  \setlength{\itemsep}{1pt}
  \setlength{\parskip}{0pt}
  \setlength{\parsep}{0pt}
}{\end{itemize}}


\begin{document}

\newcommand{\HRule}{\rule{\linewidth}{0.6mm}}
\begin{titlepage}

\begin{center}

\sffamily{
  \Large{
    \textbf{Tesis de la carrera\\
            Licenciatura en Ciencias de la Computación\\
    }
  }
}

\vspace*{2cm}

  {\Large{\textbf{\rmfamily{\sc Generador de Evaluador Estático para Gramática de Atributos Multi-plan\\}}}}
\vspace*{0.3cm}
\HRule \\
  \vspace*{0.6cm}
%   {\Huge\textbf{\maggen\\}}
{\centering \includegraphics[width=150pt, height=35pt]{maggen.png}}
%   \vspace*{0.1cm}
 
\HRule \\[0.5cm]

\vspace*{1.5cm}

\large{\textbf{Autores:}\\ \textbf{Kilmurray}, Gerardo Luis\\ \textbf{Picco}, Gonzalo Martín\\}
\vspace*{0.6cm}
{------\\}
\vspace*{0.6cm}
\large{\textbf{Director:}\\ Mg. Arroyo Marcelo Daniel\\}
\vspace*{1cm}
{------\\}
\vspace*{2cm}
 \includegraphics[width=40px,height=60.2px]{unrc.jpg}\\
\vspace*{0.5cm}
\normalsize{\textit{Departamento de Computación\\
                    Facultad de Ciencias Exactas, Físico-Químicas y Naturales\\
                    Universidad Nacional de Río Cuarto\\
                    Córdoba - Argentina}}
\end{center}
\end{titlepage}
\sloppy

\titlepage


\dominitoc

\pagenumbering{roman}

 \cleardoublepage

\section*{Acknowledgments}

Last thing to do :-)

\tableofcontents

\mainmatter

\chapter{Introducci\'on}
\label{chap:intro}
\minitoc


Desde que D. Knuth introdujo en 1966 las gramáticas de atributos (GA), estas se han utilizado ampliamente para el desarrollo de herramientas de procesamiento de lenguajes formales como compiladores, intérpretes, traductores como también para especificar la semántica de lenguajes de programación. Las gramáticas de atributos son un formalismo simple para la especificación de la semántica de lenguajes formales, como los lenguajes de programación o de especificación. Integran la modularidad que brindan las gramáticas libres de contexto y la expresividad de un lenguaje funcional.

\section{Gramática de Atributos}

En una gramática de atributos, se relaciona con cada símbolo de una gramática libre de contexto un conjunto de atributos. Cada regla o producción tiene asociados un conjunto de reglas semánticas que toman la forma de asignación a atributos de valores denotados por la aplicación de una función, la cual puede tomar como argumentos intancias de atributos pertenecientes a los símbolos que aparecen en la producción.
Las reglas semánticas inducen dependencias entre los atributos que ocurren en la producción. El orden de evaluación es implícito (si existe) y queda determinado por las dependencias entre las instancias de los atributos.
Una regla semántica se podrá evaluar cuando las instancias de los atributos que aparecen como sus argumentos estén evaluadas. Un evaluador de gramáticas de atributos debe tener en cuenta las dependencias entre las instancias de atributos para seguir un orden consistente de evaluación de los mismos.
Si una GA contiene dependencias circulares no podría ser evaluada ya que no podría encontrarse un orden de evaluación. Existen numerosas herramientas basadas en este formalismo o en alguna de sus ex-tensiones, entre las cuales podemos mencionar yacc, Yet Another Compiler-Compiler [29],desarrollado por ATT, AntLR, JavaCC, JavaCUP, ELI y muchas otras.


\section{Métodos de Evaluación}

Los métodos estáticos deben tener en cuenta todos los posibles árboles sintácticos posibles a ser generados por la gramática y calcular todas las posibles dependencias entre las instancias de los atributos. Además, se deberán detectar las posibles depen-dencias circulares, para informar la viabilidad de su evaluación.
Esto se conoce como el problema de la circularidad, el cual se ha demostrado ser intrínsecamente exponencial [20]. El problema de la circularidad ha motivado que muchos investigadores hayan rea-lizado esfuerzos en la búsqueda e identificación de familias o subgrupos de gramáticas de atributos, para las cuales puedan detectarse circularidades con algoritmos de menor complejidad (polinomial o lineal).
Estas familias imponen restricciones sobre la gramática de atributos o sobre las dependencias entre sus atributos para garantizar que una GA no sea circular, con el costo de restringir su poder expresivo. 
Las clases de familias de gramáticas de atributos que se han utilizado para el desarrollo de herramientas eficientes y que se encuentran ampliamente analizadas en la bibliografía especializada, encontramos las s-atribuidas2 , l-atribuidas, las gramáticas de atributos ordenadas (OAG) y las absolutamente no circulares (ANCAG)[2]. En 1980, Uwe Kastens[23] caracterizó las gramáticas de atributos ordenadas y propuso un método para su evaluación, denominado secuencias de visita. Estas son secuencias de operaciones que conducen el recorrido del árbol sintáctico atribuido y realizan la evaluación de las instancias de los atributos. Kastens propone un método para generar las secuencias de visita en tiempo poli-nomial para la familia OAG.
Mas recientemente, en 1999, se han propuesto nuevas familias de GA para las que se pueden implementar evaluadores eficientes basado en métodos estáticos y con un mayor poder expresivo que las utilizadas tradicionalmente[44].



\subsection{Evaluci\'on din\'amica}

bla bla

\subsection{Evaluaci\'on est\'atica}

\subsection{Evaluacion de la familia ANCAG}
\section{Secuencia de visita}
\section{Generaci\'on de evaluadores para GA bien definidas}
\section{Evaluaci\'on durante el parsing}

bla bla



\appendix

\chapter{Apéndice}
\label{chap:appendix}

\section{Especificación completa en \spirit}
\label{append:grammarspirit}

\begin{lstlisting}[numbers=left,basicstyle=\tiny,language=C++,numberstyle=\tiny, numbersep=5pt]
/**
  * Declaration of the Attribute Grammar structure
  * with the Spirit library of Boost.
  */
struct attritute_grammar: public grammar<attritute_grammar>
{
  template <typename ScannerT>
  struct definition
  {
    definition(attritute_grammar const &self)
    {
      /* Common declarations. */

      r_ident = lexeme_d[(alpha_p|'_')>>*(alnum_p|'_')]-r_reserved_word;

      r_oper = lexeme_d[(alpha_p|'_'|r_id_op)>>*(alnum_p|'_'|r_id_op)];

      r_id_op = ch_p('+')|'*'|'/'|'^'|'%'|'&'|'<'|'='|'-'|'>'|'|'|'~'|'.'|','|'?';

      r_boolean = str_p("true")|"false";

      r_char = lexeme_d[ch_p('\'')>>(anychar_p)>>ch_p('\'')];

      r_string = lexeme_d[ch_p('\"')>>r_string_lit>>ch_p('\"')];

      r_string_lit = +((anychar_p-(ch_p('\"')|"\\"|'\'' ))|r_esc_seq);

      r_esc_seq = ch_p('\\')>>
                  ( oct_p
                  | as_lower_d['x']>>hex_p
                  | (anychar_p-ch_p('\n'))
                  )
                ;

      r_reserved_word = str_p("compute")|"end"
                      | "all"
                      | "semantic domain"|"attributes" |"rules"
                      | "sort"|"op"|"function"
                      | "infix"|"prefix"|"postfix"
                      | "syn"|"inh"
                      | "left"|"right"|"non-assoc"
                      | r_cpp_reserved_words
                      ;

      r_cpp_reserved_words = r_cpp_basic_types
                           | str_p("and")|"and_eq"|"asm"|"auto"|"bitand"
                           | "bitor"|"break"|"case"|"catch"|"class"|"compl"
                           | "const"|"const_cast"|"continue"|"default"
                           | "delete"|"do"|"double"|"dynamic_cast"|"else"
                           | "enum"|"explicit"|"export"|"extern"|"false"
                           | "for"|"friend"|"goto"|"if"|"inline"|"long"
                           | "mutable"|"namespace"|"new"|"not"|"not_eq"
                           | "operator"|"or"|"or_eq"|"private"|"protected"
                           | "public"|"register"|"reinterpret_cast"|"return"
                           | "short"|"signed"|"sizeof"|"static"|"static_cast"
                           | "struct"|"switch"|"template"|"this"|"throw"|"true"
                           | "try"|"typedef"|"typeid"|"typename"|"union"
                           | "unsigned"|"using"|"virtual"|"void"|"volatile"
                           | "wchar_t"|"while"|"xor"|"xor_eq"
                           ;

      r_cpp_basic_types = str_p("bool")|"char"|"float"|"int"|"string";

      /*
       * Declaration of Semantic Domain.
       */
      r_semantic_domain = lexeme_d[str_p("semantic domain")>>space_p]>>
                          +r_bloq_sem;

      r_bloq_sem = r_decl_sort
                 | r_decl_oper[&add_operator]
                 | r_decl_func[&add_function];

      /* Declaration of Sorts. */

      r_decl_sort = lexeme_d[str_p("sort")>>space_p]>>
                    (r_ident[&create_sort][st_sorts.add]%',')>>
                    ';';

      /* Declaration of Operators. */

      r_decl_oper  = lexeme_d[str_p("op")>>space_p][&inic_func]>>
                     (r_oper_infix|r_oper_postfix|r_oper_prefix)>>
                     str_p("->")>>
                     r_sort_st[&save_image_func]>>
                    ';';

      r_oper_infix = str_p("infix")[&save_mode_op]>>
                     !r_oper_mode>>
                     r_oper[&save_name_func][st_op_infix.add]>>
                     ':'>>
                     r_sort_st[&save_domain_func]>>','>>r_sort_st[&save_domain_func];

      r_oper_postfix = str_p("postfix")[&save_mode_op]>>
                      !r_oper_mode>>
                      r_oper[&save_name_func][st_op_postfix.add]>>
                      ':'>>
                      r_sort_st[&save_domain_func];

      r_oper_prefix = !(str_p("prefix")[&save_mode_op])>>
                      !r_oper_mode>>
                      r_oper[&save_name_func][st_op_prefix.add]>>
                      ':'>>
                      r_sort_st[&save_domain_func];

      r_oper_mode  = '('>>
                      (uint_p[&save_prec_op]|'_')>>
                      ','>>
                      (r_oper_assoc[&save_assoc_op]|'_')>>
                      ')';

      r_oper_assoc = str_p("left")|"right"|"non-assoc";

      /* Declaration of Functions. */

      r_decl_func  = lexeme_d[str_p("function")>>space_p]>>
                     r_oper[&inic_func][&save_name_func][st_functions.add]>>
                     ':'>>
                     !r_dom_func>>
                     str_p("->")>>
                     r_sort_st[&save_image_func]>>
                     ';';

      r_dom_func  = r_sort_st[&save_domain_func]%',';

      /*
       * Declaration of Attributes.
       */
      r_attributes = lexeme_d[str_p("attributes")>>space_p]>>
                     +r_decl_attr[&create_attributes];

      r_decl_attr = (r_ident[&add_attribute][st_attributes.add]%',')>>
                    ':'>>
                    !(r_type_attr[&save_type_attr])>>
                    '<'>>r_sort_st[&save_sort_attr]>>'>'>>
                    lexeme_d[str_p("of")>>space_p]>>
                    (r_conj_symb |
                    (str_p("all")>>!('-'>>r_conj_symb))
                    )[&save_member_list_attr]>>
                    ';';

      r_conj_symb = '{'>>(r_ident%',')>>'}';

      r_type_attr = str_p("inh")|"syn";

      /*
       * Declaration of Rules.
       */

      r_rules = lexeme_d[str_p("rules")>>space_p]>>
                (+r_decl_rule)>>eps_p[&check_well_defined];

      r_decl_rule = r_ident[&create_new_non_terminal][&create_rule][st_non_terminal.add]>>
                    str_p("::=")>>
                    r_right_rule[&save_rule]>>
                    *(str_p("|")[&create_abbreviated_rule]>>
                    r_right_rule[&save_rule])>>
                    ';';

      r_right_rule = +( r_ident[&create_new_non_terminal][st_non_terminal.add]
                      | r_terminal[&create_new_terminal]
                      )[&save_right_side_rule]>>
                     !r_compute_eq;

      r_compute_eq = str_p("compute")>>
                     +(r_equation)>>
                     str_p("end");

      r_terminal = lexeme_d[ch_p('\'')>>r_string_lit>>ch_p('\'')];

      r_equation = r_instance[&create_equation]>>
                   '='>>
                   r_expression[&save_rvalue]>>
                   ';';

      /*
       * Expression's Grammar non ambiguos based in
       *
       *    E = T *(<op_infix> T)
       *    T = F *(<op_postfix>)
       *      |(<op_prefix>) T
       *    F = (E)
       *      |function
       *      |literal
       *      |instance
       */
      r_expression = r_expr_prime>>*(r_op_infix_st[&create_operator][&create_func_node]>>
                     r_expr_prime[&create_root_infix_node]);

      r_expr_prime = r_expr_prime_prime>>
                     *(
                      r_op_postfix_st[&create_operator][&create_func_node][&create_root_postfix_node]
                      )
                   | r_op_prefix_st[&create_operator][&create_func_node]>>
                     r_expr_prime[&create_root_prefix_node]
                   ;

      r_expr_prime_prime  = ch_p('(')[&increment_level]>>r_expression>>ch_p(')')[&decrement_level]
                          | r_function[&create_root_function_node]
                          | r_literal[&create_literal_node]
                          | r_instance[&create_instance_node]
                          ;

      /*
       * The functions accept a list of expressions.
       */
      r_function = r_function_st[&create_function][&create_func_node]>>
                   ch_p('(')[&push_mark][&increment_level]>>
                   !(r_expression%',')>>
                   ch_p(')')[&decrement_level];

      /*
       * Literals accepted: Integer and Float numbers, characters
       * and string, between signs ' and " respectively.
       */
      r_literal = longest_d[real_p|int_p][&create_lit_number]
                | r_char[&create_lit_ch]
                | r_string[&create_lit_str]
                | r_boolean[&create_bool]
                ;

      /*
       * An instance is, the symbol with the number of occurrences in
       * square brackets within the rule, with the specific attribute
       * with which it operates.
       *
       * Example: E[0].value
       */
      r_instance = lexeme_d[
                     r_non_term_st[&create_instance]>>
                     '['>>int_p[&save_index_ins]>>']'>>
                     '.'>>r_attribute_st[&save_attr_ins]
                   ];

      /*
       * Declaration of Attribute Grammar.
       */
      r_att_grammar = r_semantic_domain>>
                      r_attributes>>
                      r_rules>>end_p ;

      /*
       * Parsers based in the symbol tables.
       */
      r_sort_st       = st_sorts|r_cpp_basic_types;
      r_op_prefix_st  = st_op_prefix;
      r_op_infix_st   = st_op_infix;
      r_op_postfix_st = st_op_postfix;
      r_function_st   = st_functions;
      r_attribute_st  = st_attributes;
      r_non_term_st   = st_non_terminal;
    }
    /*
     * Symbols's Table for the elements of an Attribute Grammar.
     */
    symbols <> st_sorts;
    symbols <> st_op_prefix;
    symbols <> st_op_infix;
    symbols <> st_op_postfix;
    symbols <> st_functions;
    symbols <> st_attributes;
    symbols <> st_non_terminal;

    /*
     * Variables using in parsing time.
     */
    typedef rule<ScannerT> rule_exp;

     /* Rule in lexeme_d. */
    typedef rule<typename lexeme_scanner<ScannerT>::type> rule_lexeme;

    /* Basic rules: identifiers and basic types. */
    rule_exp r_reserved_word, r_cpp_reserved_words, r_cpp_basic_types,
             r_ident, r_oper, r_char, r_string, r_boolean;

    rule_lexeme r_id_op, r_string_lit, r_esc_seq;

    /* Semantic domain's rule: Sort, Operator and Function. */
    rule_exp r_semantic_domain, r_bloq_sem, r_decl_sort, r_decl_oper,
             r_decl_func, r_oper_assoc, r_oper_mode, r_oper_prefix,
             r_oper_infix, r_oper_postfix, r_dom_func;

    /* Atribute's rule. */
    rule_exp r_attributes, r_decl_attr, r_type_attr, r_conj_symb;

    /* Rule's rule. */
    rule_exp r_rules, r_decl_rule, r_equation, r_right_rule,
             r_terminal, r_compute_eq;

    /* Expresion's rule: Compute. Add context for type expresion. */
    rule_exp r_expression, r_expr_prime, r_expr_prime_prime, r_function,
             r_literal, r_instance;

    /* Translate for symbol table. */
    rule_exp r_sort_st, r_op_prefix_st, r_op_infix_st, r_op_postfix_st,
             r_function_st, r_attribute_st, r_non_term_st, r_sort_stable;

    /* Main rule. */
    rule_exp r_att_grammar;

    rule_exp const &start() const
    {
      return r_att_grammar;
    }
  };
};
\end{lstlisting}

\section{Ejemplo: Attribute grammar Wuu Yang}
\label{append:agwuuyang}
\scriptsize
\subsection{Pseudocodigo}
\begin{lstlisting}[escapeinside=@@, backgroundcolor=\color{white}]
     (R1)   S @$\rightarrow$@ XYZ      
                     S.s0 := X.s1 + Y.s2 + Ys3 + Zs4
                     X.i1 := Y.s3  
                     Y.i2 := X.s1
                     Y.i3 := Y.s2
     (R2)   Y @$\rightarrow$@ m        
                     Y.s2 := Y.i2
                     Y.s3 := 1
     (R3)   Y @$\rightarrow$@ n        
                     Y.s2 := 2
                     Y.s3 := Y.i3
     (R4)   X @$\rightarrow$@ m        
                     X.s1 := X.i1
     (R5)   Z @$\rightarrow$@ Y        
                     Z.s4 := Y.s3
                     Y.i2 := 3
                     Y.i3 := Y.s2
\end{lstlisting} 

\subsection{Con lenguaje de especificación}
\scriptsize
\begin{lstlisting}[numbers=left, numberstyle=\tiny, numbersep=5pt, language=cobol ]
/**
  *  \file    ag_Wuu_yang.input
  *  \brief   Attribute Grammar example.
  *  \date    15/02/2010
  *  \author  Kilmurray, Gerardo Luis <gerakilmurray@gmail.com>
  *  \author  Picco, Gonzalo Martin <gonzalopicco@gmail.com>
  */

// Block of Semantic Domain
semantic domain
  // List of Operators 
  op infix  (10, left) +: int, int -> int;

// Block of Attributes
attributes
  s0 : syn <int> of {S};
  s1 : syn <int> of {X};
  s2 : syn <int> of {Y};
  s3 : syn <int> of {Y};
  s4 : syn <int> of {Z};
  i1 : inh <int> of {X};
  i2 : inh <int> of {Y};
  i3 : inh <int> of {Y};

// Block of Rules
rules
  // R1
  S ::= X Y Z
      compute  
          S[0].s0 = X[0].s1 + Y[0].s2 + Y[0].s3 + Z[0].s4;
          X[0].i1 = Y[0].s3;
          Y[0].i2 = X[0].s1;
          Y[0].i3 = Y[0].s2;
      end;

  // R2
  Y ::= 'm'
      compute
          Y[0].s2 = Y[0].i2;
          Y[0].s3 = 1;
      end;

  // R3
  Y ::= 'n'
      compute
          Y[0].s2 = 2;
          Y[0].s3 = Y[0].i3;
      end;

  // R4
  X ::= 'm'
      compute
          X[0].s1 = X[0].i1;
      end;
  
  // R5
  Z ::= Y
      compute
          Z[0].s4 = Y[0].s3;
          Y[0].i2 = 3;
          Y[0].i3 = Y[0].s2;
      end;           
\end{lstlisting}


\bibliographystyle{ThesisStyle}
\bibliography{Thesis}

%\printnomenclature

\cleardoublepage
\begin{vcenterpage}
\noindent\rule[2pt]{\textwidth}{0.5pt}
\begin{center}
{\large\textbf{Design and Use of Numerical Anatomical Atlases for Radiotherapy\\}}
\end{center}
{\large\textbf{Abstract:}}
The main objective of this thesis is to provide radio-oncology specialists with automatic tools for delineating organs at risk of a patient undergoing a radiotherapy treatment of cerebral or head and neck tumors.
\\
To achieve this goal, we use an anatomical atlas, i.e. a representative anatomy associated to a clinical image representing it. The registration of this atlas allows to segment automatically the patient structures and to accelerate this process. Contributions in this method are presented on three axes.
\\
First, we want to obtain a registration method which is as independent as possible w.r.t. the setting of its parameters. This setting, done by the clinician, indeed needs to be minimal while guaranteeing a robust result. We therefore propose registration methods allowing to better control the obtained transformation, using outlier rejection techniques or locally affine transformations.
\\
The second axis is dedicated to the consideration of structures associated with the presence of the tumor. These structures, not present in the atlas, indeed lead to local errors in the atlas-based segmentation. We therefore propose methods to delineate these structures and take them into account in the registration.
\\
Finally, we present the construction of an anatomical atlas of the head and neck region and its evaluation on a database of patients. We show in this part the feasibility of the use of an atlas for this region, as well as a simple method to evaluate the registration methods used to build an atlas.
\\
All this research work has been implemented in a commercial software (Imago from DOSIsoft), allowing us to validate our results in clinical conditions.
\\
{\large\textbf{Keywords:}}
Atlas-based Segmentation, non rigid registration, radiotherapy, atlas creation
\\
\noindent\rule[2pt]{\textwidth}{0.5pt}
\end{vcenterpage}

\end{document}
